%!TEX root = ../chapter04.tex 
%----------------------------------------------------------------------------------------
%	
%----------------------------------------------------------------------------------------



%----------------------------------------------------------------------------------------
%	 
%----------------------------------------------------------------------------------------


\newpage 
\loadgeometry{marginlayout}  
\pagestyle{scrheadings} % Print headers again  
\KOMAoptions{
	headwidth=textwithmarginpar,
	headsepline=true,
	headsepline= 0.5pt,
	footwidth=textwithmarginpar,
}   
\onehalfspacing   


\section{Nombres premiers}
\begin{definition}
	Un nombre entier est \textbf{premiers} s'il admet \textbf{exactement} 2 diviseurs \textbf{différents} : 1 et lui même.
\end{definition}
On classe les nombres entiers $\mathbb{N}^*$ en 3 catégories :
\begin{description}
	\item[l'unité] 1
	\item[les nombres premiers] 2, 3, 5...
	\item[les nombres composés] tout le reste.
\end{description}
\begin{theorem}[Nombres premiers < 100]   sont :
	\begin{enumerate}[label=\textbullet, leftmargin=*]
		\item $2$, $3$, $5$ et $7$.
		\item tous les nombres non divisibles par 2, 3, 5 et 7.
	\end{enumerate}
\end{theorem}
\begin{example}
	$11$ , $31$ et $97$ sont des nombres premiers.
\end{example} 
\begin{theorem}[Test de primalité (admis)]  
	Si $n\in\mathbb{N}^*$ n'admet pour diviseur aucun des nombres (premiers) compris entre 2 et $\sqrt{n}$, alors $n$ est premier.
\end{theorem}
\begin{example}
	Vérifier que les nombres $149$ et $173$ sont premiers.
\end{example}

\begin{theorem}[théorème fondamental de l'arithmétique]  Tout nombre entier $n\in\mathbb{N}^*$ se décompose de manière unique (à l'ordre près) en produit de  facteurs premiers.
\end{theorem}
\begin{example}
	$31$ est un nombre premier, sa décomposition en facteurs premiers est lui-même.
	
	$1008=2^4\times 3^2 \times 7^1$. Le facteur 2 a pour multiplicité 4.
	
	$25777=149\times 173$.
\end{example} 
%----------------------------------------------------------------------------------------
%	EXERCICES
%----------------------------------------------------------------------------------------

\newpage 
\loadgeometry{widelayout}  
\pagestyle{scrheadings} % Print headers again  
\KOMAoptions{
	headwidth=textwithmarginpar,
	headsepline=true,
	headsepline= 0.5pt,
	footwidth=textwithmarginpar,
} 
\onehalfspacing


\subsection{TP Algorithmes de test de primalité} \label{chap03:exo:python:premier01}

\begin{tBox}
	\begin{enumerate}[label={\textbullet}, leftmargin=*]   
	\item l'instruction \lstinline[language=python ,style=python-code]|a//b| retourne le quotient de la division entière de $a$ par $b$.
	\item l'instruction \lstinline[language=python ,style=python-code]|a%b| retourne le reste de la division entière de $a$ par $b$.
	\end{enumerate}
\end{tBox} 
L'objectif est de comprendre le fonctionnement des différents scripts par une lecture attentive. Les programmer sur une calculatrice peut servir de vérification, mais n'est ni nécessaire ni suffisant.
\begin{exercise}\hfill  
\begin{multicols}{2} 
\begin{lstlisting}[ language=python , 
%	title=Programme A,   
linewidth=1\columnwidth,
style=python-code,    
aboveskip=0.2\topskip,
belowskip=0.2\topskip,
xleftmargin=1.2em,
tabsize=4,
%	firstnumber=1  
]
def fonction(n) :
	x = 0
	for i in range(1,n):
		if n%i == 0 :
			x = x+1
	return x
def estpremier(n) :
	p = fonction(n)
	if ....
		return True
	else :
		return False
\end{lstlisting}
\end{multicols} 
\begin{enumerate}[label={\arabic*)}, leftmargin=*]   
	\item On utilise l'instruction \lstinline[language=python ,style=python-code]|fonction(12)|. Faire tourner le script à la main, et vérifier que la valeur retournée est 5. Préciser :
		\begin{enumerate}[label={\alph*)}, leftmargin=*]   
			\item le nombre de répétitions de la boucle \lstinline|for|.
			\item pour chaque répétition \lstinline|i|, la valeur de \lstinline|x| en fin.
		\end{enumerate}
	\item Lancer à la main l'instruction \lstinline[language=python ,style=python-code]|fonction(5)| et \lstinline[language=python ,style=python-code]|fonction(6)|.
	\item On suppose que \lstinline[language=python ,style=python-code]|fonction(n)| est un entier $\geqslant 2$. 
	Corriger la ligne $3$ afin que l'instruction \lstinline[language=python ,style=python-code]|fonction(n)| retourne le nombre de diviseurs de \lstinline[language=python ,style=python-code]|fonction(n)|.
	\item Comment utiliser l'instruction \lstinline|fonction(n)| pour déterminer si \lstinline|n| est premier ?
	\item Compléter la ligne 9 afin que l'instruction \lstinline|estpremier(n)| retourne \lstinline|True| si \lstinline|n| est premier, et \lstinline|False| dans le cas contraire. 
\end{enumerate} 
\end{exercise}
\begin{exercise}\label{chap03:exo:python:premier02}\hfill  
\begin{multicols}{2} 
\begin{lstlisting}[ language=python , 
%	title=Programme A,   
linewidth=1\columnwidth,
style=python-code,    
aboveskip=0.2\topskip,
belowskip=0.2\topskip,
xleftmargin=1.2em,
tabsize=4,
%	firstnumber=1  
]
def fonction(n) :
	i = 2
	while n%i !=0 :
		i = i + 1
	return i
	
def estpremier(n) :
	p = fonction(n)
	if ......
		return True
	else :
		return False
\end{lstlisting}
\end{multicols} 
	\begin{enumerate}[label={\arabic*)}, leftmargin=*]   
		\item  Faire tourner le script à la main, et vérifier que l'instruction \lstinline|fonction(7)| retourne \lstinline[language=python]|7|.
	\item Que retourne l'instruction \lstinline|fonction(12)| ?  \lstinline|fonction(35)| ? 
	\item Que retourne l'instruction \lstinline|fonction(n)| lorsque \lstinline|n| est un nombre entier supérieur ou égal à 2 ?
	\item Comment utiliser l'instruction \lstinline|fonction(n)| pour déterminer si \lstinline|n| est premier ?
	\item Compléter la ligne 9 afin que l'instruction \lstinline|estpremier(n)| retourne \lstinline|True| si \lstinline|n| est premier, et \lstinline|False| dans le cas contraire.
	\end{enumerate} 
\end{exercise}



\begin{exercise}\label{chap03:exo:python:premier03} \hfill 

	\begin{enumerate}[label={\arabic*)}, leftmargin=*]   
	\item Expliquer pourquoi \lstinline|10101| n'est pas premier. 
	\item Comparer le nombre de boucles nécessaire à chacun des algorithmes des exercices \ref{chap03:exo:python:premier01} et \ref{chap03:exo:python:premier02} pour identifier que  \lstinline|10101| n'est pas premier. Quel algorithme est le plus rapide ?
	\item Le script ci dessous  \lstinline|estpremier()| ci-dessous est supposé retourner \lstinline|True| si la variable en entrée \lstinline|n| est un nombre premier et \lstinline|False| sinon.
	
\noindent\begin{lstlisting}[ language=python , 
	%	title=Programme A,   
	linewidth=0.95\columnwidth,
	style=python-code,    
	aboveskip=0.2\topskip,
	belowskip=0.2\topskip,
	lineskip=1.1em,
	xleftmargin=1.2em,
	tabsize=4,
	%	firstnumber=1  
	]
def estpremier(n) :
	x = 2
	while x < n**0.5 :		# rappel : $n^{\numprint{0.5}}=\sqrt{n}$, alternative x**2<n
		if n%x== 0 :		
			return ...
		else :
			x = x + 1
	return ...
\end{lstlisting}	
	On rappelle que l'exécution de la fonction s'arrête à l'exécution d'une instruction  \lstinline|return|. 
	\begin{enumerate}[label={\alph*)}, leftmargin=*]   
		\item Faire tourner le script avec l'instruction \lstinline|estpremier(15)|.\\
			Combien de fois la boucle \lstinline|while| est exécutée ?\\
			Que vaut la variable \lstinline|x| à l'arrêt de la fonction ?
		\item Mêmes questions, cette fois on utilise l'instruction \lstinline|estpremier(15)|.
		\item Compléter les lignes 5 et 8 par \lstinline|True| ou \lstinline|False|.
	\end{enumerate} 
	\item Expliquer pourquoi il n'est pas nécessaire de chercher des diviseurs au delà de $\sqrt{n}$. 
	\item Faire tourner (à la main) le scipt avec l'instruction \lstinline|estpremier(9)| puis \lstinline|estpremier(16)|. Que constatez vous ?
	\item Corriger la ligne 3, afin de remédier à l'erreur de la question précédente.
	\end{enumerate} 
\end{exercise}
\begin{exercise}[Réfuter une conjecture à l'aide d'un algorithme]\hfill 
  
Pierre de Fermat (\textborn\,1605, \textdied\,1665) pensait que tous les entiers $F_n= 2^{(2^n)}+1$ étaient des nombres premiers.  Effectivement $F_0=3$ , $F_1=5$ et $F_2=2^{(2^2)}+1=17$ sont des nombres premiers.
  
À l'aide de l'algorithme final de l'exercice \ref{chap03:exo:python:premier03} , et en complétant le script ci-dessous, trouve le plus petit entier $F_n$ qui n'est pas premier.
 
\noindent\begin{lstlisting}[ language=python , 
	%	title=Programme A,   
	linewidth=0.95\columnwidth,
	style=python-code,    
	aboveskip=0.2\topskip,
	belowskip=0.2\topskip,
	lineskip=1.1em,
		xleftmargin=1.2em,
	tabsize=4,
		firstnumber=10  
	]
n , F = 0 , 3
while ....
	n = n + 1
	F =
print(n)
\end{lstlisting} 
\end{exercise}
 

%----------------------------------------------------------------------------------------
%	
%----------------------------------------------------------------------------------------